\SourceLeaf{543}
\pdfbookmark[0]{Magnetic Needle in Copenhagen and Milan. 1834}{copenhagen_mailand_1834}
\begin{center}\itshape
Observations of Variations of the Magnetic Needle in Copenhagen and Milan\\
on 5 and 6 November 1834.
\end{center}
\begin{center}
\hrule\vspace{2mm}
\textsc{Schumacher}. Astronomische Nachrichten No.\ 276. 21 March 1835.
\vspace{2mm}\hrule\vspace{5mm}
\end{center}
Since the completion of the magnetic observatory here, variations in magnetic declination have been observed regularly, among other things, at short intervals without interruption on certain days fixed in advance. Initially the dates chosen were the same as those that Mr \textsc{von Humboldt} had arranged several years earlier. Since last spring, quite a number of astronomers and physicists have acquired instruments similar to ours, which I have described in sufficient detail elsewhere, and are taking part in the agreed observations. The very first simultaneous observations obtained in this way, on 4 and 5 May in Göttingen and Waltershausen (an estate near Schweinfurt, where Mr \textsc{Sartorius} observed with a small instrument otherwise very like ours), showed a striking degree of agreement in the way the variations repeatedly jumped back and forth, in the smaller movements as well as the larger ones. Similar results have often been obtained since then at later sessions in which Leipzig, Berlin, Braunschweig and Copenhagen took part. Some examples have been published as graphs in \textsc{Poggendorff}'s Annalen.

As participation in these agreed observations continues to spread, we can naturally expect increasingly interesting and fruitful results. I therefore repeat the announcement already made elsewhere. Since it has become so clear that observations must be made at \textit{very} short intervals, we have found it necessary to change the schedule. We have reduced the number of sessions from 8 to 6 per year and their duration from 44 to 24 hours. The current arrangement is the last Saturday of every odd-numbered month, from noon at Göttingen to noon the following day. Two additional sessions accompany each of these principal sessions, on
\SourceLeaf{544}the following Tuesday and Wednesday evenings from 8 to 10 o'clock. A more detailed account, including the method of observation, can be found in \textsc{Poggendorff}'s Annalen, vol.\ 33. [p.\ 528 of this volume.]

The most remarkable result obtained so far comes from the simultaneous observations at Copenhagen and Milan on 5 and 6 November of this year. This was a session under the earlier arrangement, and the observers at those places had yet to receive notice of the change. In Copenhagen, where a magnetic observatory very like ours has now been established under the direction of State Councillor \textsc{Oersted}, a needle of the same strength as ours was used (weighing four pounds). In Milan, Mr \textsc{Sartorius} and Dr \textsc{Listing} observed at the local astronomical observatory with the smaller needle mentioned above, assisted by Mr \textsc{Kreil}, a pupil at the observatory. I admit that even after the many earlier experiences of a similar kind, I was surprised at the degree of agreement between two places more than 150 miles apart. The accompanying graph speaks for itself. I add only a few explanations and remarks.

At first I did not know the value of the scale divisions in Copenhagen, so I drew the graph to a scale that made the anomalies appear about equally large. I achieved this by making the side of each square in the grid correspond to nine scale divisions for the Copenhagen observations and three for the Milan observations. In fact one scale division in Copenhagen amounts to $21''576$ and one in Milan to $29''341$. Measured as angles, the movements in Copenhagen were therefore about $2{,}2$ times as large as those in Milan. To infer the ratio of the forces at work, one must bear in mind that these phenomena are to be regarded only as disturbances of the horizontal force of the Earth's magnetism at the two places. The same disturbing force must produce greater changes where the horizontal force is smaller than where it is greater. From \textsc{Hansteen}'s fine map in volume 7 of the Astronomische Nachrichten\HistoricalNote{544}, I estimated the ratio of the horizontal force of the Earth's magnetism in Copenhagen and Milan as
\SourceLeaf{545}1 to $1{,}23$. On that basis, the ratio of the disturbing forces that produced the most considerable anomalies at Copenhagen and Milan on those days could be estimated at about $1{,}8$ to 1. How much better we will be able to judge such puzzling forces of nature when we have similar simultaneous observations at many places far apart.

Alongside the surprisingly close agreement in the course of the anomalies, we certainly notice differences as well. Yet what seems surprising is how small those differences are in comparison. We still know nothing of the causes of these phenomena. In view of their varied interplay, it seems unnatural to suppose that all the causes acted from a single point. Some may have been situated here and others there, and during the 44 hours forces from quite different regions may well have come into play in quite a different ratio at the two places. The curve for Milan generally looks much more jagged than the one for Copenhagen. This is readily explained by the fact that observations were made every 5 minutes at the former place and every 10 minutes at the latter. At the longer intervals some of the smaller and more rapidly changing anomalies must therefore have gone unnoticed.

These investigations hardly need anything to make them more interesting, but I believe I must point out one circumstance of particular concern to astronomers. We do not yet know whether the forces acting in these movements take a measurable time to propagate over great distances. This interesting question will doubtless be answered in
\SourceLeaf{546}the future. If the time is too small to measure, observations of such rapid rising and falling movements can serve to determine longitudes. Under favourable circumstances these determinations may hold their own even against the more precise determinations of longitude. From movements observed in Göttingen and Leipzig I have already several times been able, on that assumption, to deduce the difference in longitude correctly to within half a minute of time. Sometimes, however, the movements are so rapid that time can be determined with much greater precision. The strongest movements I have encountered so far occurred on 7 February of this year, when the needle was extremely restless all day. I observed movements of 17 scale divisions or 6 minutes of arc in one minute of time. They continued steadily for several minutes, then gradually slowed and afterwards went in the opposite direction. If such phenomena were followed carefully at two places with good instruments (giving an accuracy of a few seconds of arc even in individual observations), they could give the difference in longitude to within one second of time, provided the forces act in a time too small to measure. In any event, to clarify the matter it will clearly be important for all observers who have the means to make sure that their time determinations are always accurate.

Finally, I note that on 2 November the observers in Milan used a \textsc{Lenoir} inclinometer to find the local inclination $=63^0\,55'\,26''$.

